
% \qh{3 figures for ss, ss-v2, and ss-pro. x-asix: train datasize, y-axis: acc.}
\section{Introduction}

% With the rapid advancement of large language models (LLMs) and large vision-language models (VLMs), there is growing interest in building GUI (Graphical User Interface) agents capable of understanding natural language instructions and autonomously interacting with software interfaces. These agents aim to simulate human-like behaviors, such as clicking buttons, filling text fields, or navigating menus across various platforms including desktop environments \cite{zhang2024ufo, zhang2025ufo2}, mobile devices \cite{wang2024mobileagentautonomousmultimodalmobile}, and web applications \cite{zheng2024gpt}. To operate effectively in real-world scenarios, a GUI agent must possess two core capabilities: (i) Multimodal perception and understanding of the digital world, allowing it to interpret both visual demonstrations and linguistic cues; and (ii) Action execution, enabling it to interact with the digital world via mouse, keyboard, and touchscreens \cite{zhang2024large}.

% % Given that rendered screenshots of digital environments are information-complete and consistent in design across different platforms,
% Early efforts in developing GUI agents often relied on structured metadata such as HTML, DOM trees, or visual hierarchy documentation \cite{zhang2024ufo}. However, these structured metadata is not always available, often contains noise or inconsistencies, and varies widely in format across different platforms. To overcome these limitations, recent research has increasingly focused on visual GUI agents. These agents perceive the interface directly from screenshots, mirroring the way humans interact with digital systems \cite{xu2024aguvis}.
% Within this paradigm, a key challenge is \textit{visual grounding}: mapping the agent's linguistic plans to precise actionable regions on the screen. Most existing methods today treat visual grounding as a text-based \textit{coordinate generation} task. In this formulation, the agent outputs screen positions (\eg, "x=0.125, y=0.23") as part of natural language generation, leveraging the same mechanisms it uses to generate any other tokenized text \cite{gou2024navigating}.

% However, coordinate-generation–based approaches come with several intrinsic limitations: \textit{(i) Weak spatial-semantic alignment}: Representing actions as discrete coordinate tokens (\eg, "x=..., y=...") requires the model to implicitly learn a mapping from visual inputs to a sequence of numeric outputs through the language modeling head. This process lacks an explicit spatial inductive bias or supervision, making it inefficient and data-hungry.
% \textit{(ii) Ambiguous supervision targets}: Many GUI actions such as clicking within a button permit a range of valid coordinates. However, coordinate regression treats this as a single-point prediction task, penalizing reasonable variations and ignoring the inherent ambiguity in such interactions.
% \textit{(iii) Feature granularity mismatch}: While screen coordinates are continuous and potentially high-resolution, visual features extracted by models like Vision Transformers (ViTs)~\cite{dosovitskiy2021an} are patch-based and relatively coarse. This mismatch in granularity forces the model to approximate ine-grained, pixel-level interactions from low-resolution, patch-level features, which can hinder generalization across different screen sizes or layouts.
% \kz{do we need to mention this part in intro?}
% Though some methods \cite{lu2024omniparser} attempt to incorporate richer spatial context by having the model predict full bounding-box coordinates instead of single-point locations, these approaches still fall short of addressing the above limitations. Bounding boxes do offer more spatial information, but when predicted as raw coordinate strings (\eg, x\_min, y\_min, x\_max, y\_max), they remain disconnected from the underlying visual features without dedicated architectural mechanisms like ROI pooling \cite{sun2019roi} or spatial attention \cite{zhu2019empirical}.


% Rethinking how humans interact with digital interfaces: \textit{people do not compute exact screen positions before acting—instead, they perceive the target and act on it directly}. Inspired by this, in this paper, we propose \sys, a VLM augmented with an attention based action head, together with a grounding verifier, enabling coordinate-free visual grounding similar to how humans do.

% Specifically, \sys introduces a dedicated \texttt{[ACTOR]} token to encode the grounding context derived from visual inputs and natural language instructions.
% To align the spatial and semantic cues relevant to the intended action, the action head employs an attention mechanism that learns to attend from the \texttt{[ACTOR]} token to visual patch tokens of the GUI screenshot. Action regions are then proposed based on the resulting attention map.
% We explored multi-patch supervision for model training, where all image patches overlapping with ground-truth bounding boxes are labeled as positives, and all others are labeled as negatives. This formulation allows the model to handle supervision ambiguity gracefully and avoids penalization of alternative valid interaction points within UI elements.
% Moreover, by aligning predictions directly at the vision backbone's spatial resolution, \sys can better generalize across screen sizes and resolutions, facilitating consistent behavior in diverse digital environments.
% Since the region proposer can suggest multiple candidate regions in a single pass, we further introduce a grounding verifier to evaluate and select the most likely region for the action execution.


With the rapid advancement of large language models (LLMs) and vision-language models (VLMs), there is increasing interest in building GUI (Graphical User Interface) agents that understand natural language instructions and autonomously interact with software interfaces across platforms such as desktops \cite{zhang2024ufo, zhang2025ufo2}, mobile devices \cite{wang2024mobileagentautonomousmultimodalmobile}, and web applications \cite{zheng2024gpt}. Effective GUI agents require two core capabilities: (i) multimodal perception to interpret visual and linguistic cues, and (ii) action execution to interact with digital environments via mouse, keyboard, or touchscreen \cite{zhang2024large, wang2024large}. While early systems relied on structured metadata (\eg HTML, DOM trees, or view hierarchies) \cite{zhang2024ufo}, such data is often noisy, inconsistent, or unavailable across platforms. Recent work thus emphasizes visual GUI agents that perceive interfaces directly from rendered screenshots, akin to human users \cite{xu2024aguvis}. A central challenge in this paradigm is \emph{visual grounding}: mapping natural language plans to screen regions. Most existing methods treat this as a coordinate generation task, producing screen positions (\eg ``x=0.125, y=0.23'') through the same text generation mechanisms used by LLMs \cite{gou2024navigating}.

However, representing GUI actions through coordinate generation, where models output screen positions as text tokens (\eg, \texttt{x=...}, \texttt{y=...}) introduces several intrinsic limitations. First, \textit{spatial-semantic alignment is weak}: generating discrete coordinate tokens requires the model to implicitly map visual inputs to numeric outputs via a language modeling head, without any explicit spatial inductive bias. This process is inefficient, data-intensive, and prone to errors due to the lack of direct supervision linking visual features to action locations. Second, \textit{supervision signals are ambiguous}: many GUI actions, such as clicking within a button, allow for a range of valid target positions. However, coordinate-based methods typically treat the task as single-point prediction, penalizing all deviations—even reasonable ones—and failing to capture the natural ambiguity of human interaction. Finally, there is a \textit{granularity mismatch between vision and action space}: while coordinates are continuous and high-resolution, vision models like Vision Transformers (ViTs)~\cite{dosovitskiy2021an} operate on patch-level features. This mismatch forces the model to infer dense, pixel-level actions from coarse visual tokens, which undermines generalization to diverse screen layouts and resolutions.

Although some recent approaches~\cite{qin2025ui} attempt to enrich spatial grounding by predicting bounding boxes instead of single points, they still represent these regions as raw coordinate strings (\eg \texttt{x\_min}, \texttt{y\_min}, \texttt{x\_max}, \texttt{y\_max}) that are detached from the visual features. Without architectural components such as ROI pooling~\cite{sun2019roi} or spatial attention mechanisms~\cite{zhu2019empirical}, such methods fall short of bridging the gap between linguistic intent and grounded visual action. %\cz{Can be removed if need space.}

% \qh{What's the model's performance of locating the same element that appears at different location? if it can recognize it on singular places, this man indiate the model "know" the meaning of this element, otherwise, the model may just learn the generation location of a text string "xxxxx".}

Rethinking how humans interact with digital interfaces reveals a key insight: \textit{humans do not calculate precise screen coordinates before acting—they perceive the target element and interact with it directly}. Motivated by this observation, we propose \sys, a VLM augmented with an attention-based action head, enabling coordinate-free visual grounding that more closely mimics human behavior. Unlike prior methods that treat action grounding as a coordinate prediction task, \sys learns to attend directly to relevant visual regions without relying on numeric coordinate generation. At the core of \sys is a dedicated \texttt{<ACTOR>} token, which encodes the grounding context by jointly processing visual input and natural language instructions. An attention mechanism then learns to align this token with the most relevant GUI regions by attending over visual patch tokens from the screenshot. The resulting attention map naturally identifies actionable regions on the interface.
% \cz{better to have a high-level figure to compare with coord-based methods.}

To address the inherent ambiguity in GUI interactions, where multiple points within a UI element (\eg a button) may all be valid, \sys is trained using multi-patch supervision. All visual patches overlapping with ground-truth bounding boxes are labeled as positives, while others are treated as negatives. This supervision strategy allows the model to tolerate spatial ambiguity and reduces over-penalization of reasonable action variations. Furthermore, because \sys grounds actions directly at the vision backbone's native spatial resolution, it avoids the granularity mismatch of previous methods and generalizes more robustly across different screen sizes, resolutions, and layouts. Finally, to support decision refinement, we further enhance GUI-Actor by presenting a lightweight grounding verifier that evaluates multiple candidate regions and selects the most plausible one for action execution.

% action head: coordinate free
% multi-region prediction with verifier
% efficient training
Our contribution can be summarized as follows:
% \begin{itemize}
\begin{enumerate}[leftmargin=*, label=\arabic*.]
    \item We revisit recent coordinate generation-based approaches for visual grounding in GUI agents, identify their limitations—such as weak spatial-semantic alignment, ambiguous supervision targets, and mismatched feature granularity—and propose \sys, a novel coordinate-free method that effectively addresses these issues.
    \item We design an attention-based action head, which can generate multiple candidate regions in a single forward pass, offering flexibility for downstream modules such as search strategies.
    \item We introduce a grounding verifier to select the most likely action region among the candidates proposed from the action attention map. We show that this verifier can be easily integrated with other grounding methods for a further performance boost.
    \item Extensive experiments demonstrate that \sys outperforms the state-of-the-art methods trained on a similar scale of data across multiple GUI action grounding benchmarks, and exhibits greater robustness to unseen screen sizes and resolutions. Remarkably, the 2B version of \sys even surpasses several competing 7B models. Furthermore, by leveraging the verifier, \sys with lightweight training (\ie, freezing the backbone LLM and fine-tuning only the newly introduced $\sim$100M parameters in the action head) can effectively equip the underlying VLM with grounding capabilities without compromising its general-purpose strengths.
\end{enumerate}
%\end{itemize}

% we propose a compelling alternative for ....
% supporting multi-region grounding
